% GENERATED FILE. Edit canonical lessons, then run npm run generate:books.
% canonical-source-hash: fnv1a64:a41a90e09f7d85c8
% canonical-lessons: TE-C414-prophesaru, TE-C414-upanyasam, TE-C414-pattabhadrudu, TE-C414-upakaravetanam, TE-C414-prasnincu

\chapter{At University}
\label{ch:te-at-university}
\hlchaptermodality{414}

\begin{chapteropening}
\textbf{By the end of this chapter:} \emph{I can talk about professors, lectures, graduates and scholarships.}

Complete the last lesson of chapter 414: I can talk about professors, lectures, graduates and scholarships.
\end{chapteropening}

\section[\texorpdfstring{\te{ప్రొఫెసరు} (prophesaru)}{ప్రొఫెసరు (prophesaru)}]{\te{ప్రొఫెసరు} (prophesaru) — a professor}

\label{lesson:TE-C414-prophesaru}



\begin{quote}
Before the new one: say the Telugu for a laboratory, then the Telugu for to invent.
\end{quote}

\subsection*{You'll want to know: \te{ప్రొఫెసరు}}
\textbf{\te{ప్రొఫెసరు}} — \emph{prophesaru} — \textquotedblleft{}a professor\textquotedblright{}.

\begin{quote}
\te{మా} \te{ప్రొఫెసరు} \te{చరిత్ర} \te{చెప్తారు}.
\end{quote}

\begin{quote}
\emph{mā prophesaru caritra ceptāru.}
\end{quote}

\begin{quote}
Our professor teaches history.
\end{quote}

Plural: \te{ప్రొఫెసర్లు}, \emph{prophesarlu}.

Like \te{బస్సు}, the English loan ends in -\te{ఉ}.

\begin{cousinweb}
\textbf{\te{ప్రొఫెసరు}} (\emph{prophesaru}) is a loan from English \emph{professor}, with a final -\te{ఉ}.
\end{cousinweb}

\begin{grammarlens}[title={What you've built}]
One more word for giving reasons and opinions.
\end{grammarlens}

\subsection*{Your turn}
\begin{itemize}
\raggedright
  \item \emph{Say it:} \emph{prophesaru}
  \item \emph{Say it:} \emph{prophesaru}, once more
  \item \emph{Say it:} \emph{mā prophesaru caritra ceptāru.} — Our professor teaches history.
  \item \emph{Recall:} say the Telugu for a laboratory, then the Telugu for to invent, then say \emph{prophesaru} again
\end{itemize}

\begin{tcolorbox}[breakable,colback=teal!4,colframe=teal!35!black,title={Before you move on}]
What does \te{ప్రొఫెసరు} mean? (\textquotedblleft{}A professor\textquotedblright{}.) Say it once more.
\end{tcolorbox}
\section[\texorpdfstring{\te{ఉపన్యాసం} (upanyāsaṁ)}{ఉపన్యాసం (upanyāsaṁ)}]{\te{ఉపన్యాసం} (upanyāsaṁ) — a lecture; a speech}

\label{lesson:TE-C414-upanyasam}



\begin{quote}
Before the new one: say the Telugu for to invent, then the Telugu for a professor.
\end{quote}

\subsection*{You'll want to know: \te{ఉపన్యాసం}}
\textbf{\te{ఉపన్యాసం}} — \emph{upanyāsaṁ} — \textquotedblleft{}a lecture; a speech\textquotedblright{}.

\begin{quote}
\te{ప్రొఫెసరు} \te{ఉపన్యాసం} \te{రెండు} \te{గంటలు} \te{సాగింది}.
\end{quote}

\begin{quote}
\emph{prophesaru upanyāsaṁ reṇḍu gaṇṭalu sāgindi.}
\end{quote}

\begin{quote}
The professor’s lecture went on for two hours.
\end{quote}

Plural: \te{ఉపన్యాసాలు}, \emph{upanyāsālu}.

\te{ఉపన్యాసం} \te{ఇచ్చు}, to give a lecture or a speech.

\begin{cousinweb}
\textbf{\te{ఉపన్యాసం}} (\emph{upanyāsaṁ}) is borrowed from Sanskrit \emph{upanyāsa}, \textquotedblleft{}placing near; bringing forward, speaking of, mention\textquotedblright{} (Monier-Williams).
\end{cousinweb}

\begin{grammarlens}[title={What you've built}]
One more word for giving reasons and opinions.
\end{grammarlens}

\subsection*{Your turn}
\begin{itemize}
\raggedright
  \item \emph{Say it:} \emph{upanyāsaṁ}
  \item \emph{Say it:} \emph{upanyāsaṁ}, once more
  \item \emph{Say it:} \emph{prophesaru upanyāsaṁ reṇḍu gaṇṭalu sāgindi.} — The professor’s lecture went on for two hours.
  \item \emph{Recall:} say the Telugu for to invent, then the Telugu for a professor, then say \emph{upanyāsaṁ} again
\end{itemize}

\begin{tcolorbox}[breakable,colback=teal!4,colframe=teal!35!black,title={Before you move on}]
What does \te{ఉపన్యాసం} mean? (\textquotedblleft{}A lecture; a speech\textquotedblright{}.) Say it once more.
\end{tcolorbox}
\section[\texorpdfstring{\te{పట్టభద్రుడు} (paṭṭabhadruḍu)}{పట్టభద్రుడు (paṭṭabhadruḍu)}]{\te{పట్టభద్రుడు} (paṭṭabhadruḍu) — a graduate}

\label{lesson:TE-C414-pattabhadrudu}



\begin{quote}
Before the new one: say the Telugu for a professor, then the Telugu for a lecture.
\end{quote}

\subsection*{You'll want to know: \te{పట్టభద్రుడు}}
\textbf{\te{పట్టభద్రుడు}} — \emph{paṭṭabhadruḍu} — \textquotedblleft{}a graduate\textquotedblright{}.

\begin{quote}
\te{మా} \te{అన్న} \te{గత} \te{సంవత్సరం} \te{పట్టభద్రుడయ్యాడు}.
\end{quote}

\begin{quote}
\emph{mā anna gata saṁvatsaraṁ paṭṭabhadruḍayyāḍu.}
\end{quote}

\begin{quote}
My elder brother graduated last year.
\end{quote}

Plural: \te{పట్టభద్రులు}, \emph{paṭṭabhadrulu}.

\te{పట్టభద్రుడు} \te{అవు}, to graduate; a woman graduate is \te{పట్టభద్రురాలు}.

\begin{cousinweb}
\textbf{\te{పట్టభద్రుడు}} (\emph{paṭṭabhadruḍu}) is built from Sanskrit \emph{paṭṭa}, \textquotedblleft{}a slab, a tablet\textquotedblright{} (CDIAL 7699), and \emph{bhadra}, \textquotedblleft{}blessed, fortunate\textquotedblright{} (CDIAL 9377), with the Telugu masculine ending -\te{డు}.
\end{cousinweb}

\begin{grammarlens}[title={What you've built}]
One more word for giving reasons and opinions.
\end{grammarlens}

\subsection*{Your turn}
\begin{itemize}
\raggedright
  \item \emph{Say it:} \emph{paṭṭabhadruḍu}
  \item \emph{Say it:} \emph{paṭṭabhadruḍu}, once more
  \item \emph{Say it:} \emph{mā anna gata saṁvatsaraṁ paṭṭabhadruḍayyāḍu.} — My elder brother graduated last year.
  \item \emph{Recall:} say the Telugu for a professor, then the Telugu for a lecture, then say \emph{paṭṭabhadruḍu} again
\end{itemize}

\begin{tcolorbox}[breakable,colback=teal!4,colframe=teal!35!black,title={Before you move on}]
What does \te{పట్టభద్రుడు} mean? (\textquotedblleft{}A graduate\textquotedblright{}.) Say it once more.
\end{tcolorbox}
\section[\texorpdfstring{\te{ఉపకారవేతనం} (upakāravētanaṁ)}{ఉపకారవేతనం (upakāravētanaṁ)}]{\te{ఉపకారవేతనం} (upakāravētanaṁ) — a scholarship}

\label{lesson:TE-C414-upakaravetanam}



\begin{quote}
Before the new one: say the Telugu for a lecture, then the Telugu for a graduate.
\end{quote}

\subsection*{You'll want to know: \te{ఉపకారవేతనం}}
\textbf{\te{ఉపకారవేతనం}} — \emph{upakāravētanaṁ} — \textquotedblleft{}a scholarship\textquotedblright{}.

\begin{quote}
\te{ఆమెకు} \te{చదువు} \te{కోసం} \te{ఉపకారవేతనం} \te{వచ్చింది}.
\end{quote}

\begin{quote}
\emph{āmeku caduvu kōsaṁ upakāravētanaṁ vaccindi.}
\end{quote}

\begin{quote}
She got a scholarship for her studies.
\end{quote}

With \te{వచ్చు}: \te{ఉపకారవేతనం} \te{వచ్చింది}, a scholarship came (to her).

\begin{cousinweb}
\textbf{\te{ఉపకారవేతనం}} (\emph{upakāravētanaṁ}) is built from Sanskrit \emph{upakāra}, \textquotedblleft{}help, use\textquotedblright{} (CDIAL 2141), and \emph{vetana}, \textquotedblleft{}wages, hire\textquotedblright{} (Monier-Williams): pay given as help.
\end{cousinweb}

\begin{grammarlens}[title={What you've built}]
One more word for giving reasons and opinions.
\end{grammarlens}

\subsection*{Your turn}
\begin{itemize}
\raggedright
  \item \emph{Say it:} \emph{upakāravētanaṁ}
  \item \emph{Say it:} \emph{upakāravētanaṁ}, once more
  \item \emph{Say it:} \emph{āmeku caduvu kōsaṁ upakāravētanaṁ vaccindi.} — She got a scholarship for her studies.
  \item \emph{Recall:} say the Telugu for a lecture, then the Telugu for a graduate, then say \emph{upakāravētanaṁ} again
\end{itemize}

\begin{tcolorbox}[breakable,colback=teal!4,colframe=teal!35!black,title={Before you move on}]
What does \te{ఉపకారవేతనం} mean? (\textquotedblleft{}A scholarship\textquotedblright{}.) Say it once more.
\end{tcolorbox}
\section[\texorpdfstring{\te{ప్రశ్నించు} (praśniñcu)}{ప్రశ్నించు (praśniñcu)}]{\te{ప్రశ్నించు} (praśniñcu) — to question; to challenge}

\label{lesson:TE-C414-prasnincu}



\begin{quote}
Before the new one: say the Telugu for a graduate, then the Telugu for a scholarship.
\end{quote}

\subsection*{You'll want to know: \te{ప్రశ్నించు}}
\textbf{\te{ప్రశ్నించు}} — \emph{praśniñcu} — \textquotedblleft{}to question; to challenge\textquotedblright{}.

\begin{quote}
\te{విద్యార్థులు} \te{ఆ} \te{నియమాన్ని} \te{ప్రశ్నించారు}.
\end{quote}

\begin{quote}
\emph{vidyārthulu ā niyamānni praśniñcāru.}
\end{quote}

\begin{quote}
The students questioned that rule.
\end{quote}

Past (I …): \te{ప్రశ్నించాను}, \emph{praśniñcānu}.

The verb of \te{ప్రశ్న}, a question.

\begin{cousinweb}
\textbf{\te{ప్రశ్నించు}} (\emph{praśniñcu}) is built on Sanskrit \emph{praśna}, \textquotedblleft{}a question\textquotedblright{} (CDIAL 8818), with the Telugu verb ending -\te{ఇంచు}.
\end{cousinweb}

\begin{grammarlens}[title={What you've built}]
One more word for giving reasons and opinions.
\end{grammarlens}

\subsection*{Your turn}
\begin{itemize}
\raggedright
  \item \emph{Say it:} \emph{praśniñcu}
  \item \emph{Say it:} \emph{praśniñcu}, once more
  \item \emph{Say it:} \emph{vidyārthulu ā niyamānni praśniñcāru.} — The students questioned that rule.
  \item \emph{Recall:} say the Telugu for a graduate, then the Telugu for a scholarship, then say \emph{praśniñcu} again
\end{itemize}

\begin{tcolorbox}[breakable,colback=teal!4,colframe=teal!35!black,title={Before you move on}]
What does \te{ప్రశ్నించు} mean? (\textquotedblleft{}To question; to challenge\textquotedblright{}.) Say it once more.
\end{tcolorbox}
